\begin{abstract}
We investigate whether large language models produce detectably different text when instructed to lie compared to truthful responses. While prior work on LLM deception detection relies on access to model internals (\eg hidden-state probes), we ask a simpler question: \emph{can deception be detected from the output text alone?}
We generate 1,344 paired responses from \gptfull and \gptmini under four conditions---truthful, direct lie, roleplay lie, and sycophantic lie---applied to 96 factual statements from the \azaria dataset. We extract 15 surface-level linguistic features and compare distributions across conditions.
Our results show that lying text is markedly different: deceptive responses are longer (Cohen's $d{=}{-}1.08$), contain 7$\times$ more negation ($d{=}{-}1.23$), and use 2.8$\times$ more hedging language ($d{=}{-}0.50$). Vocabulary divergence between truthful and deceptive conditions is 3.7--4.4$\times$ the truthful self-divergence baseline. A Random Forest classifier achieves 92.8\% accuracy (AUC 0.971) at distinguishing truthful from deceptive text using only these surface features---no model internals needed. These signatures generalize across all three deception strategies and across both models tested. Our findings suggest that current LLMs do not convincingly imitate their own truthful style when lying, opening a practical path toward black-box deception detection.
\end{abstract}
